\chapter{The game and its requirements}
\label{chap-exigences}

\section{Rules}

The player commands the fleet of the \emph{Sup Galilée Marine Nationale} against a fleet of
pirates led by the computer. Each side has a grid of $10 \times 10$ cells and a fleet of ten
ships, that is twenty occupied cells (\cref{tab-flotte}).

\begin{table}[ht]
    \centering
    \begin{tabular}{@{}lccc@{}}
        \toprule
        Type of ship   & Length (cells) & Number      & Occupied cells \\
        \midrule
        Four-decker    & 4              & 1           & 4              \\
        Three-decker   & 3              & 2           & 6              \\
        Two-decker     & 2              & 3           & 6              \\
        One-decker     & 1              & 4           & 4              \\
        \midrule
        \textbf{Fleet} &                & \textbf{10} & \textbf{20}    \\
        \bottomrule
    \end{tabular}
    \caption{Composition of a fleet.}
    \label{tab-flotte}
\end{table}

\begin{description}[style=nextline, leftmargin=1.2em]
    \item[Placement] Ships are straight, horizontal or vertical. Two ships never touch,
          \emph{neither by a side nor by a corner}: each ship is surrounded by a belt of water
          (\cref{fig-non-contact}). The fleets are placed at random; the player can have their own
          redrawn as long as no shot has been fired.
    \item[Turn] The player starts. Firing means designating a cell of the opposing grid. If a
          ship is there, it is \emph{hit} and the shooter plays again; otherwise the shot is a
          \emph{miss} and the turn passes.
    \item[Sunk ship] A ship all of whose cells are hit is \emph{sunk}.
    \item[End of the game] The first side whose whole fleet is sunk has lost.
\end{description}

\begin{figure}[ht]
    \centering
    \begin{tikzpicture}[scale=0.62]
        % --------------------------------------------------- (a) legal placement
        \begin{scope}
            \fill[ecume] (1,2) rectangle (6,5);
            \plateau{7}
            \fill[marine, rounded corners=1.5pt] (2.12,3.12) rectangle (4.88,3.88);
            \fill[ocean, rounded corners=1.5pt] (6.12,0.12) rectangle (6.88,1.88);
            \fill[ocean, rounded corners=1.5pt] (0.12,6.12) rectangle (1.88,6.88);
            \node[note] at (3.5,-0.8) {(a) a ship's belt of water\\and legal neighbours};
        \end{scope}
        % ----------------------------------------------- (b) forbidden placements
        \begin{scope}[xshift=9.5cm]
            \plateau{7}
            \fill[marine, rounded corners=1.5pt] (2.12,3.12) rectangle (4.88,3.88);
            \fill[corail, rounded corners=1.5pt] (5.12,4.12) rectangle (5.88,5.88);
            \fill[corail, rounded corners=1.5pt] (1.12,2.12) rectangle (3.88,2.88);
            \draw[corail, line width=1.2pt] (4.75,3.75) -- (5.25,4.25) (4.75,4.25) -- (5.25,3.75);
            \draw[corail, line width=1.2pt] (2.75,2.75) -- (3.25,3.25) (2.75,3.25) -- (3.25,2.75);
            \node[note] at (3.5,-0.8) {(b) contact by a corner\\and contact by a side};
        \end{scope}
    \end{tikzpicture}
    \caption{The no-contact rule. In (a), the light area is forbidden to the other ships; in
        (b), two placements that are refused.}
    \label{fig-non-contact}
\end{figure}

The no-contact rule has a consequence that shapes the whole project: it makes it possible to
\emph{deduce}. A hit proves that its four diagonal neighbours are water; a sunk ship proves
that its whole belt is too. The game marks these cells with a flag, and the strength of the
artificial opponent depends first of all on the use it makes of these deductions
(\cref{chap-adversaire}).

\section{Functional requirements}

\begin{table}[ht]
    \centering
    \begin{tabularx}{\textwidth}{@{}lX@{}}
        \toprule
        \textbf{F1} & Play a complete game against the computer, by the rules above.          \\
        \textbf{F2} & Choose among three difficulty levels: easy, medium, hard.               \\
        \textbf{F3} & Place the fleets at random; redraw the player's before the first shot.  \\
        \textbf{F4} & Display the state of both grids: the player's ships, misses, hits, sunk
        ships, cells excluded by deduction.                                                   \\
        \textbf{F5} & Display the score of each side and whose turn it is.                    \\
        \textbf{F6} & Consult the rules and the presentation of the game at any time.         \\
        \textbf{F7} & At the end of a game, offer to play again, change level or quit.        \\
        \textbf{F8} & Play with the keyboard as well as the mouse; optional sound effects.    \\
        \bottomrule
    \end{tabularx}
    \caption{Functional requirements.}
    \label{tab-fonctionnelles}
\end{table}

\section{Design requirements}

These are what set version~2 apart from the first. Each of them can be verified, and
\cref{chap-qualite} shows by what means it actually is.

\begin{table}[ht]
    \centering
    \begin{tabularx}{\textwidth}{@{}lXl@{}}
        \toprule
                    & \textbf{Requirement}                                             & \textbf{Verified by} \\
        \midrule
        \textbf{C1} & The model depends neither on GTK nor on any other library.       &
        build without GTK                                                                                     \\
        \textbf{C2} & One seed replays exactly the same game, on every platform.       &
        reference file                                                                                        \\
        \textbf{C3} & The opponent cannot cheat: it does not know where the ships are. &
        construction (types)                                                                                  \\
        \textbf{C4} & A refused operation leaves the state of the game untouched.      & unit tests           \\
        \textbf{C5} & No compiler warning; no memory error.                            &
        \code{-Werror}, ASan, UBSan                                                                           \\
        \textbf{C6} & The program starts from any folder.                              &
        embedded resources                                                                                    \\
        \textbf{C7} & The interface never freezes during the computer's turn.          &
        one shot per clock tick                                                                               \\
        \bottomrule
    \end{tabularx}
    \caption{Design requirements and means of verification.}
    \label{tab-conception}
\end{table}
